\documentclass[11pt,a4paper]{article}
\usepackage[utf8]{inputenc}
\usepackage[T1]{fontenc}
\usepackage{amsmath, amssymb, amsthm}
\usepackage{hyperref}
\usepackage{graphicx}
\usepackage{geometry}
\usepackage{xcolor}
\usepackage{enumitem}

\geometry{margin=1in}

\title{\textbf{Identity-Driven Discourse Systems (IDDS) 2.1: \\ A Framework for Automated Conflict Modeling and Sovereign Moderation}}
\author{Johanna Almeida \\ \textit{Project Falcon Maintainer}}
\date{March 2026}

\begin{document}

\maketitle

\begin{abstract}
The Identity-Driven Discourse System (IDDS) provides a formal taxonomy and transition model for online communication, mapping the degradation from rational exchange to non-recoverable conflict. Version 2.1 introduces the $D_{flag}$ (Disagreement) modifier, identifying Identity Activation as a context-dependent floating multiplier rather than a fixed state. This paper details new findings on Moral Protective Framing (MPF), Adversarial Seeding, and the "Silence Bypass" phenomenon. We outline the implementation of IDDS within AI-driven "Sovereign Moderation" tools designed to identify and mitigate adversarial escalation before it reaches a terminal state.
\end{abstract}

\section{Introduction}
Modern online discourse is increasingly characterized by rapid escalation into high-conflict states. As digital identity matures through protocols like AT Protocol, the interaction between individual identity and collective discourse becomes a critical vector for system stability. The IDDS framework (version 2.0) introduces a systematic approach to categorizing these interactions and predicting their failure points.

\section{Taxonomy of Identity Layers}
At the core of IDDS is the realization that conflict intensity is directly related to the "depth" of identity layers activated during an interaction. We categorize identity into three compounding layers:

\begin{itemize}[leftmargin=1.5in]
    \item[\textbf{Personal Identity}] The most granular layer, encompassing lived experiences, self-perception, and individual history. Activation here leads to defensive posturing regarding one's character.
    \item[\textbf{Ideological Identity}] The layer of belief systems, socio-political frameworks, and core values. This includes deeply held convictions (e.g., socio-economic theories, gender essentialism) that provide a map for navigating the world.
    \item[\textbf{Group Identity}] The broadest layer, involving external affiliations such as nationality, ethnicity, or cultural belonging. When this layer is triggered, the individual acts as a proxy for the collective, often resulting in "National Defense" mechanisms.
\end{itemize}

\section{The Escalation Model}
Conflict intensity $I$ is modeled as a product of identity activation, active disagreement, and environmental multipliers. We propose the revised IDDS 2.1 escalation probability formula:

\begin{equation}
    P(E) = \sigma \left( \sum_{i \in \{P, I, G\}} \omega_i \cdot A_i \cdot D_{flag} + S \cdot \text{TopicMultiplier} \right)
\end{equation}

where:
\begin{itemize}
    \item $A_i$ represents the activation state of the Personal ($P$), Ideological ($I$), and Group ($G$) layers.
    \item $D_{flag}$ is the Disagreement flag ($1$ if active conflict is present, $0$ otherwise). Identity activation only accelerates escalation when disagreement is established.
    \item $\omega_i$ are weights assigned to each layer's sensitivity.
    \item $S$ is the inherent sensitivity of the topic multiplier.
    \item $\sigma$ is a squashing function representing the threshold of discourse collapse.
\end{itemize}

Simultaneous activation of all three layers in an adversarial context ($D_{flag} = 1$) marks a critical transition point where discourse is highly likely to fail. We distinguish between \textbf{Voluntary Disclosure} (neutral context, $D_{flag} = 0$) and \textbf{Forced Exposure} (adversarial context, $D_{flag} = 1$). 

\section{Discourse State Transition Model}
IDDS identifies six distinct states of online interaction. Version 2.1 refines \textbf{Identity Activation} as a floating modifier:

\begin{enumerate}
    \item \textbf{Neutral}: Fact-based inquiry, information sharing, or benign social bonding.
    \item \textbf{Disagreement}: Divergent viewpoints expressed without personal or ideological friction. Sets $D_{flag} = 1$.
    \item \textbf{Identity Activation (Floating Modifier)}: The pivot from the objective topic toward the self. 
    \begin{itemize}
        \item If $D_{flag} = 1$: Identity activation triggers rapid escalation.
        \item If $D_{flag} = 0$: Identity activation remains a neutral social disclosure.
    \end{itemize}
    \item \textbf{Personalization}: The focus shifts entirely to the individual's character or perceived flaws.
    \item \textbf{Ad Hominem}: Rational engagement is replaced by direct, often vitriolic attacks.
    \item \textbf{Dogpile (Non-recoverable)}: A collective state of hostility. Note that while terminal for a thread, \textbf{Transient Dogpile Groups} often persist across multiple threads.
\end{enumerate}

\section{Defensive Mechanisms and Systemic Amplifiers}
Participants under "Identity Threat" employ documented defensive mechanisms, including the newly identified \textbf{Moral Protective Framing (MPF)}:
\begin{itemize}
    \item \textbf{Victim Labeling \& Minimization}: Reframing power dynamics to invalidate opposing stances.
    \item \textbf{National Defense}: Leveraging Group Identity to shield against valid ideological critique.
    \item \textbf{Competence Attacks}: Targeting the perceived authority or intelligence of the opponent.
    \item \textbf{Whataboutism}: Deflecting challenges through perceived hypocrisy in the opponent's "out-group."
    \item \textbf{Moral Protective Framing (MPF)}: Invoking a third-party vulnerability (children, family, nation) to justify escalatory or violent behavior. MPF acts as a \textbf{transition accelerator} that masks state.
\end{itemize}

These mechanisms are amplified by platform design: \textbf{Anonymity}, \textbf{Performance Incentives}, and the \textbf{Silence Bypass} (where users redirect conflicts to their own timelines to circumvent local thread moderation/blocking).

\section{Implementation in AI-Driven Moderation}
The IDDS framework enables high-fidelity "Sovereign Moderation":
\begin{enumerate}
    \item \textbf{Escalation Detection}: Machine learning models trained on the IDDS states can classify comments in real-time.
    \item \textbf{Transition Mapping}: Identifying the "pivot" point where a user shifts from disagreement to identity activation allows for early-warning interventions.
    \item \textbf{Targeted Intervention}: Rather than simple bans, AI can flag specific mechanisms (like Whataboutism or Dogpiling) for nuanced community feedback or automated "cool-down" periods.
\end{enumerate}
\section{Identity-Driven Discourse System (IDDS) --- Extended Findings}
This section captures empirical observations derived from naturalistic probing of online discourse systems (specifically Reddit), focusing on escalation dynamics, internal drivers, and control mechanisms.

\subsection{Core Escalation Triggers}
\begin{itemize}
    \item \textbf{Gender Identity Trigger}: Activates personal + ideological identity layers. Produces rapid escalation and definitional conflict, often leading to dismissal and ad hominem behavior.
    \item \textbf{National Identity Trigger}: Activates group identity. Produces defensive responses (e.g., “Our country is not like that,” “Go live somewhere else”) as a mechanism to protect the perceived attack on the collective self.
    \item \textbf{Combined Identity Activation}: Gender + National identity trigger creates multi-layer activation, resulting in maximum escalation probability and frequently leading to a dogpile state.
\end{itemize}

\subsection{Perception Defense (Model Inertia)}
\textbf{Definition}: Tendency to preserve an existing mental model of reality despite conflicting evidence.
\begin{itemize}
    \item \textbf{Observed Behavior}: Dismissal of new information, minimization of issues, and reframing criticism as exaggeration.
    \item \textbf{Mechanism}: Cognitive dissonance avoidance and identity-linked system perception.
\end{itemize}

\subsection{Ego Closure and Meta-Analysis Mechanisms}
\begin{itemize}
    \item \textbf{Final-Word Bias (Ego Closure Mechanism)}: Tendency to seek the last response as symbolic dominance or closure. This prolongs escalation cycles and prevents natural disengagement.
    \item \textbf{Meta-Analysis Trigger}: Attempting to explain behavior during conflict increases aggression. This happens due to a perceived loss of control, identity exposure (“being analyzed”), status asymmetry, and invalidation of agency. 
    \item \textbf{Insight}: People accept system explanations, but resist being explained themselves.
\end{itemize}

\subsection{Internal Drivers of Escalation}
\begin{itemize}
    \item \textbf{Perceived Threat}: Identity or beliefs feel challenged, leading to defensive response.
    \item \textbf{Identity Fragility}: Rigid or externally validated identity that requires constant defense.
    \item \textbf{Low Introspection Capacity}: Difficulty questioning own beliefs leads to externalization (e.g., “you are wrong” instead of “maybe I’m wrong”).
    \item \textbf{Emotional Overflow}: Cognitive overload converted into attack or dismissal.
\end{itemize}

    \item \textbf{Mechanism}: Threat to self-concept (“I am competent”) and control restoration behavior.
\end{itemize}

\section{IDDS 2.1 --- New Addendums and Findings}

\subsection{Adversarial Seeding}
Adversarial seeding occurs when a post is structurally designed to provoke Identity Activation in replies. In such cases, the thread is "born escalated" ($D_{flag} = 1$ at $T=0$). Detection signals include high-identity-load content paired with open engagement prompts ("what do you think?") in a new thread context.

\subsection{De-escalation Artifacts and Positive Reinforcement}
\begin{itemize}
    \item \textbf{De-escalation Artifacts}: Humor and non-sequiturs function as local state resets, breaking reinforcement loops by introducing non-threatening stimuli.
    \item \textbf{Positive Reinforcement}: Direct validation of a participant's identity or argument can resolve the underlying threat. However, insincerity triggers the \textbf{Competence Defense}, accelerating escalation.
\end{itemize}

\subsection{Transient Dogpile Groups}
Dogpile state is often not a terminal state for the group, only for the thread. Persistent groups cycle through targets, maintaining "escalation momentum." Detection requires cross-thread actor clustering.

\section{Dataset and Labeling Schema}
Empirical validation was conducted across 16 labeled screenshots from Threads, Reddit, and WhatsApp (Portuguese/English). The automated scraper utilizes the following schema:

\begin{verbatim}
post_id, anon_user_id, text, local_state, global_state, 
d_flag, mpf_flag, topic_multiplier, language, thread_id
\end{verbatim}

\section{Open Questions for 2.1}
\begin{itemize}
    \item Formalize $\omega_i$ weight learning from labeled data.
    \item Cross-thread propagation detection at scale.
    \item Formal conditions for de-escalation transitions (Reversibility).
    \item MPF detection model vs. standard sentiment analysis.
    \item Integration into Juntos/Falcon Sovereign Moderation layers.
\end{itemize}

\section{Conclusion}
By formalizing the relationship between identity and discourse, IDDS 2.1 provides a roadmap for building resilient digital communication platforms. The addition of the $D_{flag}$ modifier and MPF detection ensures that moderation systems can distinguish between sincere identity disclosure and weaponized identity exposure, moving closer to a truly "Sovereign" moderation layer for the AT Protocol ecosystem.

\end{document}
